% c03 — source pages 22–31, questions 116–169 (idioms, phrasal verbs, modals, language functions)
\begin{questions}

\Q{116} ‘సాధారణ దృక్పథం’ అనే అర్థం వచ్చే సరైన జాతీయాన్ని \en{(Idiomatic expression)} ఎంచుకోండి.
\opts{\en{a bird’s eye view}}{\en{a lion’s eye view}}{\en{a cat’s eye view}}{\en{a owl’s eye view}}

\Q{117} ‘గుర్తించని ప్రమాదం’ అనే అర్థం వచ్చే సరైన జాతీయాన్ని \en{(Idiomatic expression)} ఎంచుకోండి.
\opts{\en{Make no bones}}{\en{Bone of connection}}{\en{Snake in the grass}}{\en{A bird’s eye view}}

\Q{118} సంచిలో నుండి \blank\ ను బయటకు వదలడం.
\eng{To let the \blank\ out of the bag.}\par
‘రహస్యాన్ని బయటపెట్టడం’ అనే అర్థం వచ్చే జాతీయంగా \en{(Idiom)} మార్చడానికి పై ఖాళీని పూరించే సరైన పదాన్ని ఎంచుకోండి.
\opts{ఎలుక \en{(rat)}}{గబ్బిలం \en{(bat)}}{చాప \en{(mat)}}{పిల్లి \en{(cat)}}

\Q{119} సరైన జాతీయాన్ని \en{(Idiomatic expression)} ఎంచుకోండి.
\opts{\en{Blood and butter}}{\en{Bread and butter}}{\en{Bread and better}}{\en{Blood and better}}

\Q{120} కింది వాక్యంలోని జాతీయాన్ని \en{(Idiom)} పూర్తి చేయడానికి సరైన పదాన్ని ఎంచుకోండి.\par
ఒక వైరస్ వ్యాపిస్తోంది. ప్రజలు \blank\ లా రాలిపోతున్నారు.
\eng{There is a virus going round. People are dropping like \blank.}
\opts{కోళ్లు \en{(chickens)}}{ఈగలు \en{(flies)}}{ఎలుకలు \en{(mice)}}{చీమలు \en{(ants)}}

\Q{121} కింది సందర్భానికి తగిన సరైన జాతీయాన్ని \en{(Idiom)} ఎంచుకోండి.\par
మీ బిడ్డను చూసి మీరు ఎంతో గర్వపడుతున్నారు.
\eng{You’re so proud of your child.}
\opts{వాడు నా కంటి పాప. \en{(He’s the apple of my eye.)}}{వాడు చూడటానికి వికారంగా ఉన్నాడు. \en{(He looks a sight.)}}{వాడు నా కళ్లకు మాత్రమే. \en{(He’s for my eyes only.)}}{వాడు నా అబ్బాయి. \en{(He’s my boy.)}}

\Q{122} సరైన జాతీయాన్ని \en{(Idiomatic expression)} ఎంచుకోండి.
\opts{\en{tone of connection}}{\en{bone of connection}}{\en{bone of contention}}{\en{tongue of contention}}

\Q{123} ‘చాలా వేగంగా’ అనే అర్థం వచ్చే అర్థవంతమైన జాతీయంగా \en{(Idiomatic expression)} మార్చడానికి ఖాళీని పూరించండి.
\eng{In \blank\ and bounds.}
\opts{ఒడి / చుట్టు \en{(laps)}}{గెంతులు \en{(leaps)}}{చప్పట్లు \en{(claps)}}{పట్టీలు \en{(bands)}}

\Q{124} ‘ఒకరిని తీవ్రంగా బాధపెట్టడం’ అనే అర్థం వచ్చే సరైన జాతీయాన్ని \en{(Idiomatic expression)} ఎంచుకోండి.
\opts{\en{make somebody’s heart}}{\en{break somebody’s heart}}{\en{hurt somebody’s heart}}{\en{break somebody’s leg}}

\Q{125} చింతించకండి! సాయంత్రం 5 గంటలకల్లా \tw{ఏదో ఒక విధంగా}{by hook or by crook} డబ్బు సిద్ధం చేస్తాం.\par
గీత గీసిన జాతీయానికి \en{(Idiomatic expression)} సరైన అర్థాన్ని ఎంచుకోండి.
\opts{ఏదో ఒక పద్ధతిలో \en{(by any method)}}{చాలా గర్వంగా \en{(very proud)}}{చాలా సులభం \en{(very simple)}}{అత్యంత చెడ్డ మార్గంలో \en{(by the worst way)}}

\Q{126} నువ్వు \tw{సరిగ్గా సరైన మాట చెప్పావు}{hit the nail on the head}.\par
గీత గీసిన జాతీయానికి \en{(Idiom)} అర్థాన్ని ఎంచుకోండి.
\opts{మేకును దెబ్బతీయడం \en{(damaging a nail)}}{ఒకరి తల పగలగొట్టడం \en{(breaking someone’s head)}}{రహస్యాలు బయటపెట్టడం \en{(disclosing the secrets)}}{సరిగ్గా సరైనదాన్ని చెప్పడం లేదా చేయడం \en{(said or done exactly the right thing)}}

\Q{127} నాకు దాన్ని వదిలేయాలని \en{(giving it up)} అనిపిస్తోంది.\par
\en{‘give up’} అనే క్రియా పదబంధం \en{(Phrasal verb)} అర్థాన్ని ఎంచుకోండి.
\opts{దేనినైనా కొనసాగించడం \en{(continue something)}}{దేనినైనా తిరస్కరించడం \en{(reject something)}}{దేనినైనా అనుమతించడం \en{(allow something)}}{ఏదైనా చేయడం ఆపివేయడం \en{(to stop doing something)}}

\Q{128} రాష్ట్రపతి ప్రతి సంవత్సరం శౌర్య పురస్కారాలను \tw{అందజేస్తారు}{gives away}.\par
\en{‘give away’} అనే క్రియా పదబంధం \en{(Phrasal verb)} అర్థాన్ని ఎంచుకోండి.
\opts{ఏదైనా కార్యక్రమాన్ని నిర్వహించడం \en{(to organize some event)}}{ఒకరిని ప్రశంసించడం \en{(to appreciate someone)}}{ఏదైనా ప్రకటించడం \en{(to declare something)}}{ఏదైనా బహూకరించడం \en{(to present something)}}

\Q{129} ప్రధానోపాధ్యాయులు ఆ సమస్యను \tw{పరిశీలిస్తారు}{look into}.\par
\en{‘look into’} అనే పదబంధానికి \en{(Phrase)} సరైన అర్థాన్ని ఎంచుకోండి.
\opts{వినియోగించడం \en{(consume)}}{తదేకంగా చూడటం \en{(stare)}}{పరిశీలించడం \en{(examine)}}{బహిర్గతం చేయడం \en{(disclose)}}

\Q{130} కూడలి వచ్చే వరకు \tw{ముందుకు సాగండి}{Carry on}, తర్వాత ఎడమ వైపు తిరగండి.\par
\en{‘carry on’} అనే క్రియా పదబంధానికి \en{(Phrasal verb)} సరైన అర్థాన్ని ఎంచుకోండి.
\opts{కదులుతూనే ఉండటం \en{(to continue moving)}}{కదలడం ఆపడం \en{(to stop moving)}}{కదలికను కాసేపు ఆపడం \en{(to pause moving)}}{దర్యాప్తు చేయడం \en{(to investigate)}}

\Q{131} విమానం ఒక గంట ఆలస్యంగా \tw{గాల్లోకి లేచింది}{took off}.\par
\en{‘took off’} అనే క్రియా పదబంధానికి \en{(Phrasal verb)} సరైన అర్థాన్ని ఎంచుకోండి.
\opts{ఎగరడం ప్రారంభించింది \en{(began to fly)}}{కిందికి దిగింది \en{(landed)}}{వాయిదా పడింది \en{(postponed)}}{చేరుకోవడం \en{(arrive at)}}

\Q{132} రైలు ఆలస్యంగా \tw{వచ్చి చేరింది}{got in}.\par
\en{‘got in’} అనే క్రియా పదబంధానికి \en{(Phrasal verb)} సరైన అర్థాన్ని ఎంచుకోండి.
\opts{బయలుదేరింది \en{(left)}}{చేరుకుంది \en{(arrived)}}{ప్రయాణించింది \en{(travelled)}}{ప్రారంభమైంది \en{(started)}}

\Q{133} పని ఒత్తిడి అతన్ని \tw{చికాకు పెట్టడం}{get to him} మొదలుపెట్టింది.\par
\en{‘get to him’} అనే క్రియా పదబంధానికి \en{(Phrasal verb)} సరైన అర్థాన్ని ఎంచుకోండి.
\opts{అతన్ని వినోదపరచడం \en{(to ammuse him)}}{అతనికి చికాకు కలిగించడం \en{(to annoy him)}}{అతన్ని ప్రోత్సహించడం \en{(to encourage him)}}{అతన్ని ఆనందపరచడం \en{(to delight him)}}

\Q{134} నేను ఊపిరి తీసుకునే వరకు ఒక్క నిమిషం \tw{ఆగండి}{Wait}.\par
\en{‘wait’} అనే పదానికి సమానమైన అర్థం గల సరైన క్రియా పదబంధాన్ని \en{(Phrasal verb)} ఎంచుకోండి.
\opts{ఆలస్యం చేయడం \en{(hold up)}}{సుదీర్ఘంగా మాట్లాడటం \en{(hold forth)}}{ఆగడం \en{(hold on)}}{వాయిదా వేయడం \en{(hold off)}}

\Q{135} మేము విహారయాత్ర ముగించేంత సేపు వర్షం \tw{కురవకుండా ఆగింది}{held off}.\par
\en{‘held off’} అనే పదబంధానికి \en{(Phrase)} సరైన అర్థాన్ని ఎంచుకోండి.
\opts{మొదలైంది \en{(started)}}{మొదలు కాలేదు \en{(did not start)}}{భారీగా మారింది \en{(became heavy)}}{ప్రారంభమైంది \en{(began)}}

\Q{136} వెబ్‌సైట్‌లో తెరిచే సమయాలను నువ్వు \tw{చూసి తెలుసుకోగలవా}{look up}?\par
\en{‘look up’} అనే క్రియా పదబంధం \en{(Phrasal verb)} అర్థాన్ని ఎంచుకోండి.
\opts{సమాచారం కోసం వెతకడం \en{(to look for information)}}{మరింత దిగజారడం \en{(to become worse)}}{వెబ్‌సైట్‌లోకి ప్రవేశించడం \en{(To enter the website)}}{వెబ్‌సైట్ నుండి నిష్క్రమించడం \en{(To leave the website)}}

\Q{137} అతని వైఖరిలో మార్పుకు ఏది \tw{కారణమైంది}{brought about}?\par
\en{‘brought about’} అనే పదబంధానికి \en{(Phrase)} సరైన అర్థాన్ని ఎంచుకోండి.
\opts{పరిమితం చేసింది \en{(ristricted)}}{నిరాకరించింది \en{(denied)}}{ఆనందపరిచింది \en{(delighted)}}{కలిగించింది \en{(caused)}}

\Q{138} మా అన్ని కంప్యూటర్ల ధరలను \tw{తగ్గించాలని}{reduce} మేము లక్ష్యంగా పెట్టుకున్నాం.\par
\en{‘reduce’} అనే పదానికి సమానమైన సరైన క్రియా పదబంధాన్ని \en{(Phrasal verb)} ఎంచుకోండి.
\opts{ఉత్పత్తి చేయడం \en{(bring forth)}}{పెంచి పోషించడం \en{(bring up)}}{కలిగించడం \en{(bring about)}}{తగ్గించడం \en{(bring down)}}

\Q{139} ఆ ఫోటోలు ఎన్నో ఆహ్లాదకరమైన జ్ఞాపకాలను \tw{గుర్తుకు తెచ్చాయి}{brought back}.\par
\en{‘brought back’} అనే పదబంధానికి \en{(Phrase)} సరైన అర్థాన్ని ఎంచుకోండి.
\opts{జ్ఞాపకాలను తిరస్కరించాయి \en{(rejected memories)}}{జ్ఞాపకాలను గుర్తుకు తెచ్చాయి \en{(made remember memories)}}{జ్ఞాపకాలను కోల్పోయాయి \en{(lost memories)}}{జ్ఞాపకాలు మాయమయ్యాయి \en{(vanished memories)}}

\Q{140} రైలు ఈ స్టేషన్‌లో \tw{కొద్దిసేపు ఆగుతుంది}{calls at}.\par
\en{‘calls at’} అనే పదబంధానికి \en{(Phrase)} సరైన అర్థాన్ని ఎంచుకోండి.
\opts{కొద్దిసేపు ఆగుతుంది \en{(stops for short time)}}{ఇక్కడ బయలుదేరుతుంది \en{(starts at)}}{వెళ్ళిపోతుంది \en{(leaves)}}{చేరుకుంటుంది \en{(arrives at)}}

\Q{141} ఈ పరిస్థితికి తక్షణ చర్య \tw{అవసరం}{needs}.\par
\en{‘needs’} అనే అర్థం వచ్చే తగిన పదబంధాన్ని \en{(Phrase)} ఎంచుకోండి.
\opts{అవసరమవడం \en{(calls for)}}{రద్దు చేయడం \en{(calls off)}}{బిగ్గరగా పిలవడం \en{(calls out)}}{కొద్దిసేపు ఆగడం \en{(call at)}}

\Q{142} ఓటింగ్ నిర్వహించాలని ఆమె ప్రభుత్వాన్ని \tw{కోరింది}{called on}.\par
\en{‘called on’} అనే పదబంధానికి \en{(Phrase)} సరైన అర్థాన్ని ఎంచుకోండి.
\opts{డిమాండ్ చేసింది \en{(demanded)}}{అంకితం చేసింది \en{(dedicated)}}{వివాదం చేసింది \en{(disputed)}}{నష్టపరిచింది \en{(damaged)}}

\Q{143} వాతావరణం బాగోలేనందున ఆట \tw{రద్దు చేయబడింది}{cancelled}.\par
\en{‘cancelled’} అనే పదానికి సరైన క్రియా పదబంధాన్ని \en{(Phrasal verb)} ఎంచుకోండి.
\opts{కోరారు \en{(called on)}}{రద్దు చేశారు \en{(called off)}}{బిగ్గరగా పిలిచారు \en{(called out)}}{ఒకరి పేరు పెట్టారు \en{(called after)}}

\Q{144} సముద్రపు వాసన ఆమెకు బాల్య జ్ఞాపకాలను \tw{గుర్తుకు తెచ్చింది}{called up}.\par
\en{‘called up’} అనే పదబంధానికి \en{(Phrase)} సరైన అర్థాన్ని ఎంచుకోండి.
\opts{జ్ఞాపకాలను తిరిగి గుర్తుకు తెచ్చింది \en{(brought the memories back)}}{జ్ఞాపకాలను కోల్పోయింది \en{(lost the memories)}}{ఆనందపరిచింది \en{(delighted)}}{మర్చిపోయింది \en{(forgot)}}

\Q{145} ఆమె ఐదుగురు పిల్లలను \tw{పెంచి పెద్ద చేసింది}{brought up}.\par
\en{‘brought up’} అనే పదబంధానికి \en{(Phrase)} సరైన అర్థాన్ని ఎంచుకోండి.
\opts{పెంచింది \en{(raised)}}{పైకి లేచింది \en{(rised)}}{విడుదల చేసింది \en{(released)}}{నియంత్రించింది \en{(restricted)}}

\Q{146} అతను లేనప్పుడు అతని వ్యవహారాలను నేను \tw{చూసుకుంటాను}{look after}.\par
\en{‘look after’} అనే క్రియా పదబంధం \en{(Phrasal verb)} అర్థాన్ని ఎంచుకోండి.
\opts{పరీక్షించడం \en{(to examine)}}{మళ్ళీ చూడటం \en{(to look again)}}{ఆశించడం \en{(to hope for)}}{జాగ్రత్తగా చూసుకోవడం \en{(to take care of)}}

\Q{147} మేము తాళం చెవుల కోసం \tw{వెతుకుతున్నాం}{looking for}.\par
\en{‘looking for’} అనే పదబంధం \en{(Phrase)} అర్థాన్ని ఎంచుకోండి.
\opts{ఆశించడం \en{(hoping for)}}{చూడటం \en{(seeing)}}{తదేకంగా చూడటం \en{(staring)}}{వెతకడం \en{(searching)}}

\Q{148} నేను లోపలికి రావచ్చా?
\eng{May I come in?}\par
ఇచ్చిన వాక్యం యొక్క భాషా విధిని \en{(Language function)} ఎంచుకోండి.
\opts{అనుమతి ఇవ్వడం \en{(Giving permission)}}{బాధ్యత \en{(Obligation)}}{అనుమతి కోరడం \en{(Seeking permission)}}{సూచన \en{(Suggestion)}}

\Q{149} బాధ్యతను \en{(Obligation)} తెలిపేందుకు ఉపయోగించే సహాయక క్రియను \en{(Modal verb)} ఎంచుకోండి.\par
నువ్వు క్రమం తప్పకుండా \blank.
\eng{You \blank\ be regular.}
\opts{\en{will}}{\en{can}}{\en{may}}{\en{must}}

\Q{150} నువ్వు చాలా బాగున్నావు!
\eng{You’re looking good!}\par
పై వాక్యం \blank\ ను తెలియజేస్తుంది.
\opts{ప్రశంస \en{(Compliment)}}{వ్యాఖ్య \en{(Comment)}}{నిబద్ధత \en{(Commitment)}}{అభినందనలు \en{(Congratulations)}}

\Q{151} మీరు పరీక్షలో ఉత్తీర్ణులైనప్పుడు అందరూ మిమ్మల్ని ఎలా అభినందిస్తారు?
\eng{When you passed an examination, how would everyone greet you?}
\opts{అభినందనలు! \en{(Congratulations!)}}{చాలా ధన్యవాదాలు \en{(Thank you very much)}}{వచ్చేసారి అదృష్టం కలిసి రావాలి! \en{(Next time better luck!)}}{మీ అందరినీ కలవడం ఆనందంగా ఉంది. \en{(Nice meeting you all.)}}

\Q{152} సరైన సహాయక క్రియను \en{(Modal verb)} ఉపయోగించి కింది సంభాషణను పూర్తి చేయండి.\par
పండితుడు: “\blank\ నీకు ఏమైనా వినిపిస్తోందా?”\par
భోజనశాల యజమాని: “అవును అయ్యా, వినిపిస్తోంది.”
\eng{The scholar: “\blank\ you hear anything?”}
\eng{The eating house keeper: “Yes sir, I can”.}
\opts{\en{May}}{\en{Can}}{\en{Are}}{\en{Don’t}}

\Q{153} ఖాళీలను పూరించడానికి సరైన సహాయక క్రియలను \en{(Modal verbs)} ఎంచుకోండి.\par
నేను నా సంచిని మర్చిపోయాను. అది ఎక్కడ ఉండవచ్చో నువ్వు ఊహించగలవా?
\eng{I forgot my bag. \blank\ you guess where it \blank\ be?}
\opts{\en{Can, may}}{\en{May, can}}{\en{May, might}}{\en{Would, would}}

\Q{154} ఈ రోజు వర్షం \blank.
\eng{It \blank\ rain today.}\par
సంభావ్యతను \en{(Possibility)} వ్యక్తపరచడానికి ఉపయోగించే సరైన సహాయక క్రియను \en{(Modal verb)} ఎంచుకోండి.
\opts{\en{shall}}{\en{would}}{\en{may}}{\en{will}}

\Q{155} మీకు ఒక కప్పు కాఫీ తీసుకురానా?
\eng{Shall I get you a cup of coffee?}\par
సరైన ఐచ్ఛికాన్ని ఎంచుకోండి.
\opts{సూచన \en{(Suggestion)}}{హెచ్చరిక \en{(Warning)}}{అనుమతి \en{(Permission)}}{సహాయ ప్రతిపాదన \en{(Offer)}}

\Q{156} కింది సందర్భంలో అనుమతిని నిరాకరించడానికి సరైన సమాధానాన్ని ఎంచుకోండి.\par
\en{A.} నేను మీ సైకిల్ వాడుకోవచ్చా?\par
\en{B.} \blank.
\eng{A. May I use your bicycle?}
\eng{B. \blank.}
\opts{\en{Yes, you may.}}{\en{No, you may not.}}{\en{Yes, you may not.}}{\en{No, you may.}}

\Q{157} “వర్షం ఆగే వరకు నువ్వు నాతో ఉండవచ్చు.”
\eng{“You can stay with me till it stops raining”.}\par
పై వాక్యంలోని భాషా విధిని \en{(Language function)} గుర్తించండి.
\opts{అనుమతి తీసుకోవడం \en{(Taking Permission)}}{అనుమతి కోరడం \en{(Seeking permission)}}{సామర్థ్యం \en{(Ability)}}{అనుమతి ఇవ్వడం \en{(Giving permission)}}

\Q{158} కింది సంభాషణను పూర్తి చేయడానికి సరైన సహాయక క్రియలను \en{(Modal verbs)} ఎంచుకోండి.\par
విద్యార్థి: నేను నీళ్లు తాగడానికి వెళ్ళవచ్చా?\par
ఉపాధ్యాయులు: అవును, వెళ్ళవచ్చు.
\eng{Student: \blank\ I go to drink water?}
\eng{Teacher: Yes, you \blank.}
\opts{\en{Can’t, can}}{\en{May, may not}}{\en{Can, can’t}}{\en{May, may}}

\Q{159} మీ బంధువు \en{(cousin)} ఎడతెరిపి లేకుండా దగ్గుతున్నారు. ఆరోగ్యం జాగ్రత్తగా చూసుకోమని వారికి సలహా ఇవ్వండి.\par
వారికి సలహా ఇచ్చే సరైన వాక్యాన్ని ఎంచుకోండి.
\opts{\en{You would take care of your health.}}{\en{You should take care of your health.}}{\en{You should to take care of your health.}}{\en{You should take care of his/her health.}}

\Q{160} సంభావ్యతను \en{(Possibility)} వ్యక్తపరచడానికి ఉపయోగించే సరైన సహాయక క్రియను \en{(Modal verb)} ఎంచుకోండి.\par
నేను రేపు ఢిల్లీకి \blank.
\eng{I \blank\ go to Delhi tomorrow.}
\opts{\en{will}}{\en{should}}{\en{must}}{\en{may}}

\Q{161} \en{A:} \blank\ దయచేసి ఈ ఉత్తరాన్ని పోస్ట్ చేస్తారా?\par
\en{B:} తప్పకుండా.
\eng{A: \blank\ you please post this letter?}
\eng{B: Certainly.}\par
ఖాళీని పూరించడానికి తగిన సహాయక క్రియను \en{(Modal verb)} ఎంచుకోండి.
\opts{\en{Should}}{\en{Must}}{\en{Could}}{\en{Do}}

\Q{162} దేవుడు నిన్ను దీవించుగాక!
\eng{May God bless you!}\par
ఇచ్చిన వాక్యం యొక్క భాషా విధిని \en{(Language function)} ఎంచుకోండి.
\opts{సంభావ్యత \en{(Possibility)}}{భవిష్యవాణి \en{(Prediction)}}{ఆకాంక్ష \en{(Wish)}}{సూచన \en{(Suggestion)}}

\Q{163} డాక్టర్‌ను సంప్రదించమని మీ స్నేహితుడికి సలహా ఇవ్వడానికి ఉపయోగించే సరైన సహాయక క్రియను \en{(Modal verb)} ఎంచుకోండి.\par
నువ్వు డాక్టర్‌ను \blank.
\eng{You \blank\ consult a doctor.}
\opts{\en{should}}{\en{might}}{\en{need}}{\en{can}}

\Q{164} ఖాళీని పూరించడానికి సరైన పదాన్ని / పదబంధాన్ని ఎంచుకోండి.\par
మనం మన స్నేహితులను \blank. వాళ్ళు గతసారి మనల్ని ఆహ్వానించారు.
\eng{We \blank\ invite our friends. They invited us last time.}
\opts{\en{has to}}{\en{had to}}{\en{need}}{\en{have to}}

\Q{165} ‘సామర్థ్యాన్ని’ \en{(Ability)} తెలియజేసే సహాయక క్రియను \en{(Modal verb)} ఎంచుకోండి.\par
ఆమె తమిళం \blank.
\eng{She \blank\ speak Tamil.}
\opts{\en{shall}}{\en{will}}{\en{should}}{\en{can}}

\Q{166} మీ సామాను నేను మోయనా?
\eng{Shall I carry your luggage?}\par
ఇచ్చిన వాక్యం యొక్క భాషా విధిని \en{(Language function)} ఎంచుకోండి.
\opts{బాధ్యత \en{(Obligation)}}{సహాయ ప్రతిపాదన \en{(Offer)}}{అనుమతి \en{(Permission)}}{అభ్యర్థన \en{(Request)}}

\Q{167} మనం నడకకు వెళ్దామా?
\eng{Shall we go for a walk?}\par
ఇచ్చిన వాక్యం యొక్క భాషా విధిని \en{(Language function)} ఎంచుకోండి.
\opts{అభ్యర్థన \en{(Request)}}{బాధ్యత \en{(Obligation)}}{సూచన \en{(Suggestion)}}{సహాయ ప్రతిపాదన \en{(Offer)}}

\Q{168} మీ పుస్తకాన్ని నాకు అరువుగా ఇవ్వగలరా?
\eng{Could you lend me your book?}\par
ఇచ్చిన వాక్యం యొక్క భాషా విధిని \en{(Language function)} ఎంచుకోండి.
\opts{అభ్యర్థన \en{(Request)}}{అనుమతి \en{(Permission)}}{సూచన \en{(Suggestion)}}{నిర్బంధం \en{(Compulsion)}}

\Q{169} వ్యాకరణపరంగా సరైన వాక్యాన్ని గుర్తించండి.
\opts{\en{The players have going to the playground.}}{\en{The players has going to the playground.}}{\en{The players had going to the playground.}}{\en{The players are going to the playground.}}

\end{questions}
